\documentclass[10pt,a4paper]{amsart}
\usepackage[margin=19mm]{geometry}
\usepackage{amsmath,amssymb,mathtools}
\usepackage[scr=rsfs,cal=euler]{mathalfa}
\usepackage{xfrac}
\usepackage[unicode,hidelinks,bookmarks=false]{hyperref}
\newcommand{\ie}{i.e.\ }
\newcommand{\cf}{cf.\ }
\newcommand{\resp}{respectively\ }
\newcommand{\Subsection}{\subsection}
\providecommand{\nobreakdash}{\nobreak\hbox{-}}
\setlength{\emergencystretch}{2em}
\multlinegap0pt
\usepackage{tikz}
\usetikzlibrary{cd}
\usepackage{amssymb}
\usepackage{xcolor}
\newtheorem{lemma}{Lemma}
\DeclareMathOperator{\Spec}{Spec}
\DeclareMathOperator{\et}{et}
\DeclareMathOperator{\geom}{geom}
\DeclareMathOperator{\piet}{\pi_1^{\et}}
\DeclareMathOperator{\pietgeo}{\pi_1^{\et, \geom}}
\title{Finiteness for \'{E}tale Fundamental Groups of N\'{e}ron Models}
\author{Frank Lu}
\date{Source: arXiv:2607.00232v1 (CC BY 4.0)}
\begin{document}
\maketitle

\noindent\emph{Source and licence.} Finiteness for \'{E}tale Fundamental Groups of N\'{e}ron Models. Version arXiv:2607.00232v1 (CC BY 4.0), distributed by arXiv under the Creative Commons Attribution 4.0 International licence, http://creativecommons.org/licenses/by/4.0/, as recorded in the arXiv licence metadata for this exact version and shown on the abstract page https://arxiv.org/abs/2607.00232v1. Full licence notice: \texttt{licenses/arxiv-2607.00232-CC-BY-4.0.txt}.\par\smallskip

\noindent\textit{Editorial scope.} Counted unit: the article's lemma geometric (source0.tex lines 155--157). The article's complete original proof of it is delivered with it (source0.tex lines 158--172, one \texttt{proof} environment of 15 lines). No mathematics is written, repaired or re-derived: the statement and the proof are the article's own lines, in the article's own notation, and no retained line is substituted or re-flowed.  No further source-local context is needed: every cross-reference of every retained line resolves inside the two delivered spans.  Nothing was omitted from any retained span, so the package carries no omission record.  No retained line contains a citation, so the wrapper carries no bibliography and no bibliography entry.  The retained lines contain the article's own diagram code, which the wrapper typesets with the standard tikz packages; no image file is delivered and the package declares no asset dependency.  Editorial formatting only, in addition: an AMS \texttt{amsart} preamble that restates the article's own declarations and macros used by the retained lines, namely the environments \texttt{lemma} on separate counters, together with the standard packages those lines need.  The retained environments print in this document as Lemma 1 in order of appearance, whereas the article prints the same items with the numbering of its own full document.  The complete native source of the article as distributed in the arXiv e-print for this exact version is bundled unchanged as \texttt{sources/204156/source0.tex}.
\begin{lemma}\label{lem: generic surjects onto integral for geometric}
Let $\mathcal{X} \rightarrow \Spec \mathcal{O}_K$ be a separated, surjective, and finite type morphism from an integral normal scheme. Suppose its generic fiber $X \rightarrow \Spec K$ is geometrically connected, and $\mathcal{X}$ has a $\mathcal{O}_K$-point. Then, we have a map $\piet(X_{\overline{K}}) \rightarrow \pietgeo(\mathcal{X})$ induced from $X_{\overline{K}} \rightarrow \mathcal{X},$ which furthermore is a surjection.
\end{lemma}

\begin{proof}
The fact that this map exists follows because $\piet(X_{\overline{K}}) \rightarrow \piet(X) \rightarrow G_K \rightarrow \piet(\Spec \mathcal{O}_K)$ is zero, and so the map $\piet(X_{\overline{K}}) \rightarrow \piet(\mathcal{X})$ lands in the kernel of $\piet(\mathcal{X}) \rightarrow \piet(\Spec \mathcal{O}_K).$
\par To show surjectivity, suppose that $g \in \pietgeo(\mathcal{X});$ we know that we can lift this to some element $\tilde{g}$ of $\piet(X)$ by the normality assumption. Let $\overline{g}$ be its image in $G_K;$ we know that $\overline{g}$ lies in the kernel of $G_K \rightarrow \piet(\Spec \mathcal{O}_K).$ By assumption, we have a $\Spec \mathcal{O}_K$ point $x: \Spec \mathcal{O}_K \rightarrow \mathcal{X}$ (giving us a rational point $x_K: \Spec K \rightarrow X$) Using the morphism $x_K,$ we can map $\overline{g}$ back to some element $(x_K)_*\overline{g} \in \piet(X).$ Since $\overline{g}$ is annihilated under the map to $\piet(\Spec \mathcal{O}_K),$ using this commutative diagram
% https://q.uiver.app/#q=WzAsNCxbMCwwLCJHX0siXSxbMCwyLCJcXHBpZXQoXFxTcGVjIFxcbWF0aGNhbHtPfV9LKSJdLFsyLDIsIlxccGlldChcXG1hdGhjYWx7WH0pIl0sWzIsMCwiXFxwaWV0KFgpIl0sWzAsMywiKHhfSylfKiJdLFszLDJdLFswLDFdLFsxLDIsInhfKiIsMl1d
\[\begin{tikzcd}
	{G_K} && {\piet(X)} \\
	\\
	{\piet(\Spec \mathcal{O}_K)} && {\piet(\mathcal{X})}
	\arrow["{(x_K)_*}", from=1-1, to=1-3]
	\arrow[from=1-1, to=3-1]
	\arrow[from=1-3, to=3-3]
	\arrow["{x_*}"', from=3-1, to=3-3]
\end{tikzcd}\]
we know that $(x_K)_*\overline{g}$ lies in the kernel of the map $\piet(X) \rightarrow \piet(\mathcal{X}).$ The element $\tilde{g} ((x_K)_*\overline{g})^{-1}$ is therefore an element of $\piet(X_{\overline{K}})$ which maps to $g,$ proving surjectivity.
\end{proof}

\end{document}
